% =============================================================================
\section{Introduction}\label{sec:intro}
% =============================================================================

The capacity to integrate long-range context is a central objective in modern sequence modeling. In regulatory genomics, distal enhancers, chromatin interactions, and gene-scale architecture couple sequence elements separated by tens to hundreds of kilobases, motivating models with large input windows. Sequence-to-function models have progressed from kilobase-scale convolutional architectures~\citep{Zhou2015DeepSEA,Kelley2016Basset,Kelley2018Basenji,Kelley2020Basenji2} to transformer-based models with explicit long-range attention~\citep{Avsec2021Enformer,Linder2025Borzoi} and sub-quadratic foundation models accepting up to $10^6$ nucleotides~\citep{Nguyen2023HyenaDNA,Nguyen2024Evo,Schiff2024Caduceus,Brixi2025Evo2}. In natural language processing (NLP), analogous progress has extended context windows from thousands to millions of tokens~\citep{Vaswani2017Transformer}. Yet across both domains, a persistent gap separates \emph{nominal context length}---the maximum input length a model architecture accepts---from the \emph{effective context} that is functionally utilized for a given prediction.

\paragraph{The nominal--effective gap.}
In NLP, the RULER benchmark~\citep{Hsieh2024RULER} demonstrated that most open-source long-context models effectively utilize less than half of their claimed context window. ``Lost in the middle'' analyses~\citep{Liu2024LostInMiddle} revealed a U-shaped performance curve: models underweight information placed at intermediate positions. In genomics, \citet{Karollus2023Limited} showed that Enformer's predictions are dominated by promoter-proximal sequence despite a 196~kb receptive field, and AlphaGenome acknowledges prediction degradation for genes more than 100~kb from the input center~\citep{Benegas2026AlphaGenome}. These observations raise a fundamental question: \emph{given a sequence model with nominal context $L$, what is the effective radius of context that materially contributes to its outputs?}

\paragraph{Why existing approaches are insufficient.}
Current methods for probing context utilization are either too coarse or too narrow to serve as a formal ECL measure. Needle-in-a-Haystack (NIAH) tests~\citep{Hsieh2024RULER} are task-specific and binary (pass/fail). Gradient-based attributions~\citep{Sundararajan2017IG} and attention-map analyses are observational rather than causal, and attention weight magnitude does not reliably indicate functional dependence~\citep{Kendiukhov2026MechInterp}. In silico saturation mutagenesis (ISM)~\citep{Nair2022fastISM} provides per-position attribution scores but does not aggregate them into a formal context-length measure with statistical guarantees. What is missing is a definition that is (i) model-agnostic, applicable to convolutional, transformer, state-space, and hybrid architectures; (ii) statistically estimable with quantified uncertainty; and (iii) equipped for formal model comparison.

\paragraph{Contributions.}
This paper addresses the gap with five main contributions:
\begin{enumerate}[leftmargin=2em]
\item \textbf{Formal definitions.} We propose perturbation--variance ECL and a family of companion quantities---the Effective Context Profile (ECP), Effective Context Dimension (ECD), directional ECL, interaction ECL, and Context Decay Spectroscopy (CDS)---that together provide a multi-faceted characterization of context utilization (\cref{sec:framework}).
\item \textbf{Theory.} We establish monotonicity, invariance, and exponential-decay ECL bounds; connect ECL to Sobol total-effect indices; prove minimax estimation rates; formalize model comparison via Blackwell ordering; and derive architectural locality bounds for convolutional, attention, and SSM architectures (\cref{sec:theory}).
\item \textbf{Estimation.} We design implementable estimators with Bernstein concentration bounds, antithetic variance reduction, importance sampling, and bootstrap confidence intervals, and prove their statistical consistency (\cref{sec:estimation}).
\item \textbf{Novel methodological ideas.} We introduce the genomic Needle-in-a-Haystack (gNIAH) protocol, Context Decay Spectroscopy, ECL-as-training-diagnostic, and cross-model ECL transfer validation (\cref{sec:framework,sec:experiments}).
\item \textbf{Experimental protocols.} We propose a comprehensive empirical program spanning ten experiments across six genomic models (Enformer, Borzoi, HyenaDNA, Caduceus, Evo~2, DNABERT-2), four perturbation types, three locus classes, and biological ground-truth validation against known long-range regulatory interactions (\cref{sec:experiments}).
\end{enumerate}

\paragraph{Why embeddings?}
Embeddings are the interface between pretraining and downstream use, enabling transfer learning, retrieval, and mechanistic interpretation. Context utilization should therefore be defined at the embedding level, not only through final task accuracy. We focus on a complementary question to representation similarity measures such as CKA~\citep{Kornblith2019CKA}: \emph{which input positions influence the embedding, and how far does that influence extend?}
